\chapter{Что читать дальше}
\sloppy

Здесь собраны источники, на которых держится этот текст, и книги для углубления. Выходные данные проверяйте по
DOI или каталогу перед цитированием.

\subsection*{Материаловедение и физика твёрдого тела в целом}
\begin{itemize}
\item Askeland D.R., Fulay P.P., Wright W.J. \emph{The Science and Engineering of Materials}. 6th ed., Cengage, 2010 ---
  обзорный учебник; по нему сверялся охват частей I, II и IV этого текста (строение, диаграммы, механические свойства).
\item Callister W.D., Rethwisch D.G. \emph{Materials Science and Engineering: An Introduction}. Wiley (любое из поздних изданий).
\item Гуляев А.П. «Металловедение». М.: Металлургия; Новиков И.И. «Теория термической обработки металлов».
\item Kittel C. \emph{Introduction to Solid State Physics}. 8th ed., Wiley, 2005 (есть русские переводы ранних изданий);
  Ashcroft N.W., Mermin N.D. \emph{Solid State Physics}. 1976 (рус. пер.: «Физика твёрдого тела», Мир, 1979).
\end{itemize}

\subsection*{Кристаллография, дифракция, микроскопия}
\begin{itemize}
\item Cullity B.D., Stock S.R. \emph{Elements of X-Ray Diffraction}. 3rd ed., Prentice Hall, 2001.
\item Горелик С.С., Скаков Ю.А., Расторгуев Л.Н. «Рентгенографический и электронно-оптический анализ». М.: МИСиС.
\item Kelly A., Knowles K.M. \emph{Crystallography and Crystal Defects}. 2nd ed., Wiley, 2012.
\item Goldstein J.I. и др. \emph{Scanning Electron Microscopy and X-Ray Microanalysis}. 4th ed., Springer, 2018;
  Williams D.B., Carter C.B. \emph{Transmission Electron Microscopy}. 2nd ed., Springer, 2009.
\item Schwartz A.J. и др. (ред.) \emph{Electron Backscatter Diffraction in Materials Science}. 2nd ed., Springer, 2009.
\end{itemize}

\subsection*{Механика и разрушение}
\begin{itemize}
\item Dieter G.E. \emph{Mechanical Metallurgy}. McGraw-Hill, 1986; Золоторевский В.С. «Механические свойства металлов».
\item Anderson T.L. \emph{Fracture Mechanics: Fundamentals and Applications}. 4th ed., CRC Press, 2017.
\item Humphreys F.J., Hatherly M. \emph{Recrystallization and Related Annealing Phenomena}. 2nd ed., Elsevier, 2004.
\item Tenckhoff E. \emph{Deformation Mechanisms, Texture, and Anisotropy in Zirconium and Zircaloy}. ASTM STP 966, 1988 ---
  скольжение, двойникование и текстура именно в цирконии.
\end{itemize}

\subsection*{Фундамент}
\begin{itemize}
\item Porter D.A., Easterling K.E., Sherif M.Y. \emph{Phase Transformations in Metals and Alloys}. 3rd ed., CRC Press, 2009 ---
  главы 1--5 покрывают главы~\ref{ch:gibbs}--\ref{ch:diagrams}, \ref{ch:diffusion}, \ref{ch:defects} и \ref{ch:nucleation} этого текста подробнее.
\item Christian J.W. \emph{The Theory of Transformations in Metals and Alloys}. Pergamon, 2002 (русский перевод: «Теория
  превращений в металлах и сплавах», Мир).
\item Hillert M. \emph{Phase Equilibria, Phase Diagrams and Phase Transformations}. Cambridge, 2008.
\item Shewmon P. \emph{Diffusion in Solids}; Mehrer H. \emph{Diffusion in Solids}. Springer, 2007.
\item Hull D., Bacon D.J. \emph{Introduction to Dislocations}. 5th ed., 2011; Kelly A., Knowles K.M.
  \emph{Crystallography and Crystal Defects}. 2012.
\item Штремель М.А. «Прочность сплавов», ч.~I--II.
\end{itemize}

\subsection*{Упругость включений и микромеханика}
\begin{itemize}
\item Eshelby J.D. The determination of the elastic field of an ellipsoidal inclusion, and related problems.
  \emph{Proc. R. Soc. A} 241 (1957) 376.
\item Mura T. \emph{Micromechanics of Defects in Solids}. 1987.
\item Хачатурян А.Г. «Теория фазовых превращений и структура твёрдых растворов». Наука, 1974
  (англ.: \emph{Theory of Structural Transformations in Solids}, 1983).
\item Moulinec H., Suquet P. \emph{Comput. Methods Appl. Mech. Eng.} 157 (1998) 69 --- решатель на FFT.
\end{itemize}

\subsection*{Пластичность и текстура}
\begin{itemize}
\item Hill R. \emph{The Mathematical Theory of Plasticity}. Oxford, 1950.
\item Kocks U.F., Tom\'e C.N., Wenk H.-R. \emph{Texture and Anisotropy}. Cambridge, 1998.
\item Engler O., Randle V. \emph{Introduction to Texture Analysis}. 2nd ed., CRC Press, 2010; Bunge H.-J.
  \emph{Texture Analysis in Materials Science}. Butterworths, 1982.
\item Bachmann F., Hielscher R., Schaeben H. Texture analysis with MTEX. \emph{Solid State Phenomena} 160 (2010) 63.
\item Kearns J.J. WAPD-TM-472 (1965); \emph{J. Nucl. Mater.} 299 (2001) 171 --- факторы Кернса.
\end{itemize}

\subsection*{Водород и гидриды циркония}
\begin{itemize}
\item Puls M.P. \emph{The Effect of Hydrogen and Hydrides on the Integrity of Zirconium Alloy Components}. Springer, 2012.
\item Motta A.T. и др. Hydrogen in zirconium alloys: A review. \emph{J. Nucl. Mater.} 518 (2019) 440.
\item Bair J., Asle Zaeem M., Tonks M. A review on hydride precipitation in zirconium alloys. \emph{J. Nucl. Mater.} 466 (2015) 12.
\item Zuzek E., Abriata J.P., San-Martin A., Manchester F.D. The H--Zr system. \emph{Bull. Alloy Phase Diagrams} 11 (1990) 385.
\item Kearns J.J. \emph{J. Nucl. Mater.} 22 (1967) 292 --- растворимость; 43 (1972) 330 --- диффузия.
\item Carpenter G.J.C. \emph{J. Nucl. Mater.} 48 (1973) 264 --- несоответствие гидрида; Westlake D.G. \emph{J. Nucl. Mater.} 26 (1968) 208 --- плоскости залегания.
\item Puls M.P. \emph{Acta Metall.} 29 (1981) 1961 --- гистерезис растворимости.
\item Займовский А.С., Никулина А.В., Решетников Н.Г. «Циркониевые сплавы в ядерной энергетике». 1994.
\item Qin W. и др. \emph{Acta Mater.} 59 (2011) 7010; Wang S., Giuliani F., Britton T.B. \emph{Acta Mater.} 169 (2019) 76 ---
  межзёренные гидриды.
\item Vizca\'{\i}no P. и др. \emph{J. Nucl. Mater.} 447 (2014) 82 --- выпадение под нагрузкой.
\item Плясов А.А. и др. \emph{Physics of Atomic Nuclei} 86(12) (2023); Ядерная физика и инжиниринг 14(1) (2023) 12 ---
  переориентация и растворимость в Э635.
\end{itemize}
